%%%%%%%%%%%%%%%%%%%%%%%%%%%%%%%%%%%%%%%%%%%%%%%%%%%%%%%%%%%%%%%%%%%%%%%%%%%%%%%%%%%%%%%%%%%%%%%%%%%%%%
%
%   Filename: abstract.tex 
%
%   Description: This file will contain your abstract.
%                 
%%%%%%%%%%%%%%%%%%%%%%%%%%%%%%%%%%%%%%%%%%%%%%%%%%%%%%%%%%%%%%%%%%%%%%%%%%%%%%%%%%%%%%%%%%%%%%%%%%%%%%

\renewcommand{\abstractname}{Abstract}
\begin{abstract}
% Consists of 150 to 250 words in a \textcolor{blue}{\textbf{single paragraph}}, that provides the reader with a summary of your research/study and what you intend to do. It should be informative enough to serve as a substitute for reading the thesis document itself. The Abstract should contain the following sentences:
% \begin{itemize}
%    \item \textbf{Motivation}. One to two sentence(s) describing the research domain or problem area.
%    \item \textbf{Problem}. One to two sentence(s) describing the research challenge or opportunity to be addressed in your study
%    \item \textbf{Contribution}. One to two sentence(s) describing your intended contribution or how you plan to address the research question 
%    \item \textbf{Methodology}. One to two sentence(s) summarizing the approach or how you will carry out your study
%    \item \textbf{Significant Results or Findings}. One to two sentence(s) describing your findings and result (This can be skipped during the proposal stage) 
% \end{itemize}

The COVID-19 pandemic has underscored the importance of cycling as an alternative mode of transportation, prompting global cities to expand their cycling infrastructure. In the Philippines, the rise in bicycle use has highlighted the need for tools that can effectively gather and utilize cyclist data to support infrastructure development focused on accessibility and safety. This study presents WeCyclePH, a crowdsourcing platform that allows cyclists to log rides, annotate hazards, and contribute media for building a cyclist-generated dataset. The methodology involved public deployment of the platform to gather ride data, GPS-based annotations, and media uploads, combined with the development of an instance segmentation computer vision model to automatically detect road hazards from contributed videos. Usability and user experience were evaluated using the User Experience Questionnaire (UEQ) and semi-structured interviews, with qualitative feedback analyzed via thematic analysis. Key results indicated high hedonic quality scores (stimulation and novelty) and positive reception toward automation features, quick-report functionality, and GPS-based annotation. Interviews emphasized the importance of safety, contribution ease, and visibility of community data as engagement drivers. The resulting dataset provides a spatially explicit, multi-source view of cycling conditions, enabling both qualitative and quantitative analyses. By unifying subjective cyclist perspectives with objective machine detections at a segment level, this approach offers a richer evidence base for evaluating infrastructure, identifying hazards, and supporting data-driven improvements to urban cycling networks.

 % The COVID-19 pandemic has underscored the importance of cycling as an alternative mode of transportation, prompting global cities to expand their cycling infrastructure. In the Philippines, the rise in bicycle use has highlighted the need for tools that can effectively gather and utilize cyclist data to support infrastructure development focused on accessibility and safety. However, current methods often lack real-time, user-centric data that reflect the true experiences of cyclists. This research addressed this opportunity by developing a crowdsourcing platform designed to collect, integrate, and compile data from cyclists, including video footage, geospatial information, and user feedback. By building a cyclist-generated dataset, this study aims to support data-driven urban planning, prioritizing safer and more accessible routes. The findings demonstrated the potential of crowdsourced data to inform and improve cycling infrastructure for safer, more inclusive transportation solutions.

% The study will implement and test the crowdsourcing tool, focusing on its usability and engagement levels among cyclists throughout Luzon to ensure effective large-scale data collection. 
% he study will provide a comprehensive dataset that can be used by urban planners, policymakers, and researchers to inform sustainable and cyclist-friendly infrastructure projects. 

% The COVID-19 pandemic has underscored the importance of cycling as a safe and sustainable mode of transportation, prompting global cities to expand their cycling infrastructure. In Metro Manila, the rise in bicycle use has highlighted the need for a thorough assessment of bike lane accessibility and safety. Traditional evaluation methods often overlook cyclists' real-world experiences, which are crucial for understanding the effectiveness of bike lanes. This research aims to address this gap by developing a crowdsourcing tool to gather cyclists' feedback on bike lane conditions, safety, and usability. By integrating this user-generated data with conventional assessment methods, the study will provide a comprehensive evaluation of Metro Manila's bike lane network. The findings will inform urban planners and policymakers, contributing to the development of safer and more accessible cycling infrastructure. This mixed-methods approach will combine quantitative data, such as GPS tracking, with qualitative feedback from cyclists to create a detailed and user-centric assessment. Ultimately, the study seeks to enhance the cycling experience and promote active transportation in Metro Manila.



%
%  Do not put citations or quotes in the abstract.
%

% Keywords can be found at \url{http://www.acm.org/about/class/class/2012?pageIndex=0}.  Click the 
% link ``HTML'' in the paragraph that starts with ''The \textbf{full CCS classification tree}...''.

\begin{flushleft}
\begin{tabular}{lp{4.25in}}
\hspace{-0.5em}\textbf{Keywords:}\hspace{0.25em} & Crowdsourcing, Cyclist Experience, Bike Lane Accessibility and Safety
\end{tabular}
\end{flushleft}
\end{abstract}
